\documentclass[]{article}

\title{JAXifying IMP modeling pipeline}
\author{Sree Ganesh Balasubramani}
\begin{document}

\maketitle
\section{IMP}
In IMP, a ``particle'' is the universal base unit, but it does not necessarily
correspond to a physical atom or a fragment of a molecule. A ``particle'' can be 
\begin{itemize}
    \item A physical residue
    \item A structural container (a chain, a molecule, a state)
    \item A restraint or a bond
    \item A rigid body reference frame
\end{itemize}
IMP allocates a large, flat, internal table to store properties. Because a model has 
hierarchies, stereochemistry restraints, and system structures, IMP allocates a lot of 
particles indices corresponding to these.  
To obtain particles that IMP creates in a model 
\begin{verbatim}
    IMP.get_particles(model, model.get_particle_indexes())
\end{verbatim}
\end{document}